\documentclass[12pt,a4paper]{article}

% Packages
\usepackage[utf8]{inputenc}
\usepackage[english]{babel}
\usepackage{amsmath}
\usepackage{amssymb}
\usepackage{graphicx}
\usepackage{hyperref}
\usepackage{geometry}
\usepackage{booktabs}
\usepackage[normalem]{ulem} % adds \sout without changing \emph behavior

% Page layout
\geometry{margin=1in}

% Document information
\title{Hierarchical Spatio-Temporal State-Space Model for Mosquito Surveillance in Cienfuegos}
\author{Rita Verstraeten}
\date{\today}

\begin{document}

\maketitle

\section{Introduction}

This document describes two complementary modelling approaches for analysing entomological surveillance data from Cienfuegos, both aimed at estimating the true mosquito presence probability $p_{bt}$ while correcting for biases introduced by mixed surveillance regimes.

\begin{description}
\item[Part~I --- Bayesian hierarchical model] A full probabilistic model implemented in Stan via \texttt{cmdstanr}. It jointly estimates ecological effects, reactive surveillance bias parameters ($\delta_0$, $\delta_1$), and spatio-temporal random effects, propagating uncertainty through all components.
\item[Part~II --- Frequentist GLMM] A maximum-likelihood binomial GLMM implemented in \texttt{glmmTMB}. It is computationally fast and used to obtain data-driven estimates of covariate directions and magnitudes, which then inform the prior distributions of the Bayesian model.
\end{description}

\section{Data Issues and Surveillance Context}

\subsection{Entomological House-Visits}

The surveillance system consists of \textbf{routine inspections} (every household, monthly), quality control re-inspections of a random subset of houses, and \textbf{reactive inspections} triggered by dengue cases (the affected household and, in some instances, nearby houses). A key difficulty is that routine and reactive inspections are not labeled in the data, absences are not recorded, and reactive effort inflates observed positivity during outbreaks.

\section{Objective}

The goal is to estimate $p_{bt} = \Pr(\text{mosquito immature stages are present in block } b \text{ at time } t)$, the true spatio-temporal mosquito presence probability, while correcting for observations that arise from two regimes: systematic surveillance (baseline, covering all households) and reactive surveillance (additional inspections triggered by dengue cases). The central challenge is that reactive surveillance inflates observed mosquito positivity during outbreaks, and the data do not label observations by surveillance type.

\part*{Part~I: Bayesian Hierarchical Model}
\addcontentsline{toc}{part}{Part~I: Bayesian Hierarchical Model}

\section{Model Framework}

\subsection{Observations}

For each block $b$ and time $t$, the dataset contains $y_{bt}$ (the total number of mosquito-positive inspection events), $N_b^{HH}$ (the number of households in the block, or universo), $C_{bt}$ (the number of dengue cases), and $X_{bt}$ (environmental covariates such as NDVI and meteorology). Operationally, every household is inspected once per period through systematic surveillance, additional inspections occur only when $C_{bt} > 0$, and reactive inspections are not separately recorded.

\subsection{Total Number of Inspection Events}

The total number of inspection events is defined as:
\begin{equation}
n_{bt} = N_b^{HH} + \kappa C_{bt},
\end{equation}
where $N_b^{HH}$ is known and fixed (universo), $C_{bt}$ is observed, and $\kappa$ is a scaling factor fixed at 2 in the current implementation. This encodes the assumption that each dengue case triggers approximately $\kappa$ additional mosquito inspections. We assume that each dengue case originates from a different household on average, so reactive inspections target distinct locations within the block.

\subsection{Latent Ecological Process}

The true mosquito presence probability is modeled as a latent spatio-temporal process:
\begin{equation}
\text{logit}(p_{bt}) = \alpha + X_{bt}\beta + u_b + v_t,
\end{equation}
where $\alpha$ is the baseline intercept, $X_{bt}$ represents environmental covariates (average temperature, relative humidity, total precipitation, NDVI), $\beta$ are regression coefficients, $u_b \sim \text{Normal}(0, \sigma_u)$ is a block-level i.i.d.\ random intercept capturing stable spatial heterogeneity not explained by covariates, and $v_t$ is a temporal random effect following an AR(1) process. Crucially, dengue cases do not enter this equation. This equation defines the quantity of interest $p_{bt}$.

\subsection{Reactive Surveillance Targeting Bias}

\begin{equation}
\text{logit}(p_{bt}^R) = \text{logit}(p_{bt}) + \delta_0 + \delta_1 \log(C_{bt}) \quad \text{for } C_{bt} > 0,
\end{equation}
where $\delta_0$ represents baseline targeting bias (reactive vs systematic), and $\delta_1$ captures increased targeting as outbreak intensity increases. This affects only the observation process, not ecology.

\subsection{Combining Systematic and Reactive Surveillance}

Each inspection event belongs to one of two types: systematic (with probability $p_{bt}$) or reactive (with probability $p_{bt}^R$). The fraction of inspections that are reactive is:
\begin{equation}
\omega_{bt} = \frac{\kappa C_{bt}}{N_b^{HH} + \kappa C_{bt}},
\end{equation}
which is implied by effort, not estimated.

\subsection{Effective Observation Probability}

The probability that a randomly selected inspection event is positive is:
\begin{equation}
\pi_{bt} = \begin{cases}
p_{bt} & \text{if } C_{bt} = 0 \\
(1 - \omega_{bt})p_{bt} + \omega_{bt}p_{bt}^R & \text{if } C_{bt} > 0
\end{cases}.
\end{equation}
This is a mixture of unbiased and biased observation processes.

\subsection{Observation Model}

\begin{equation}
y_{bt} \sim \text{Binomial}(n_{bt}, \pi_{bt}).
\end{equation}

\textbf{Note on implementation:} In the Stan model, $\kappa$ is fixed (not estimated), so $n_{bt}$ is fully determined by the data. The Binomial likelihood is used directly without approximation.

\section{Extensions}

\subsection{Distributed Lag Model}

To consider lag effects of predictors, we model the effect of covariate $X_{bt}$ as:
\begin{equation}
\text{Effect of } X_{bt} = \sum_{\ell=0}^{L} w_{k,\ell} X_{bt-\ell},
\end{equation}
where $w_{k,\ell}$ is the lag weight (estimated) for covariate $k$ at lag $\ell$, and $L$ is the maximum ecological memory. To enforce smoothness in the lag structure and prevent overfitting, we impose a random walk prior on the lag weights:
\begin{align}
w_{k,1} &\sim \text{Normal}(0, 0.5), \\
w_{k,\ell} &\sim \text{Normal}(w_{k,\ell-1}, \sigma_w[k]) \quad \text{for } \ell = 2, \ldots, L+1, \\
\sigma_w[k] &\sim \text{Exponential}(2).
\end{align}
This prior encourages neighboring lag weights to be similar, resulting in smooth lag trajectories that are more interpretable and less prone to overfitting. The original equation becomes:
\begin{equation}
\text{logit}(p_{bt}) = \alpha + \underbrace{\sum_{k=1}^{K_L} \sum_{\ell=0}^{L} w_{k,\ell} X_{k,b,t-\ell}}_{\text{lagged covariates}} +  \underbrace{\sum_{k=1}^{K_U} w_{k} X_{k,b}}_{\text{static spatial covariates}} + s_b + u_t,
\end{equation}
where $k$ indexes the predictor, and $K_L$ represents the predictors that require lagging, $K_U$ represent the static spatial predictors without lagging, $w_{k,\ell}$ are lag weights following the random walk prior, $s_b$ are spatial random effects, and $u_t$ are temporal random effects (AR).

\subsection{Temporal Autocorrelation}

The temporal random effect is modeled with autoregression:
\begin{equation}
v_t = \rho v_{t-1} + \epsilon_t,
\end{equation}
where $\rho$ is the autoregressive parameter and $\epsilon_t$ is white noise.

\subsubsection{Sources of Autocorrelation}

Several factors drive autocorrelation in $p_{bt}$. Slowly varying environments (NDVI, urban form), seasonal climate, block-level structural risk, and imperfect control effectiveness all tend to increase temporal correlation. Conversely, larvae removal upon detection, targeted interventions during outbreaks, and reactive surveillance activities tend to reset states and reduce autocorrelation. Despite these opposing forces, we expect environmental suitability to dominate: Aedes recolonization is rapid, so monthly aggregation smoothens the effects of removal, and blocks remain ecologically suitable even after larvae are removed (high-risk remains high-risk).



\section{Implementation}

\subsection{Dataset Preparation}

The analysis requires a dataset structured with one row per time--block combination. Each row must contain columns for $C_{bt}$ (dengue cases), all environmental covariates $X_{bt}$, the number of houses and larvae in block $b$ at time $t$, and $N_{HH}$ (the number of houses in that block, held constant across time).

\subsection{Priors}
The information on priors selection is available in Table 1. 

\begin{table}[h]
\centering
\small
\begin{tabular}{p{1.5cm}p{3.0cm}p{3.0cm}p{3.0cm}p{5cm}}
\toprule
Parameter & Role in Model & Suggested Prior & Interpretation & Reasoning \\
\midrule
\textbf{$\alpha$} & Baseline log-odds of mosquito positivity & Normal(-4.5, 1.2) & Baseline mosquito presence when covariates = 0 & Shifted way down for rare events. On logit scale, -4.5 $\approx$ 0.011 probability, reflecting low baseline positivity in absence of favorable conditions. \\
\addlinespace
\textbf{$\delta_0$} & Baseline reactive targeting boost & Normal(0.3, 0.4) & Log-odds increase under reactive inspection & OR ($e^{0.3}$) $\approx$ 1.35 typical. Slightly reduced baseline targeting bias to stabilize initialization. \\
\addlinespace
\textbf{$\delta_1$} & Sensitivity to outbreak intensity & Normal(0, 0.2) & Log-linear response to log($C_{bt}$) & Reduced to stabilize initialization; encourages small positive effect but allows no effect; avoids explosive targeting. \\
\addlinespace
\textbf{$\rho$} & AR(1) temporal persistence & Normal(0, 0.2), constrained to $(-1, 1)$ & Temporal autocorrelation & Tighter prior regularizes toward independence. Allows positive (mosquito persistence) and negative (removal after detection) autocorrelation. The constraint is enforced by the parameter declaration in Stan (\texttt{real<lower=-1,upper=1>}), not by prior truncation. \\
\addlinespace
\textbf{$u_{block\_raw}$} & Standardized spatial effects & Normal(0, 1) & Non-centered parameter & Computational device only; ensures stable sampling. \\
\addlinespace
\textbf{$v_{time\_raw}$} & AR(1) \textit{innovation} terms & Normal(0, 1) & Standardized shocks & Computational innovation term in AR(1). The random shock entering the system at time $t$. \\
\addlinespace
\textbf{$\sigma_v$} & SD of temporal innovations & Exponential(3) & Size of month-to-month ecological shocks & Stricter prior due to AR(1) structure; prevents volatility while allowing moderate temporal variation. \\
\addlinespace
\textbf{$\sigma_u$} & SD of spatial random effects & Exponential(2) & Between-block heterogeneity & Tighter prior (prevents extreme variance during initialization). Shrinks block effects to prevent them from absorbing ecological signal. \\
\addlinespace
\textbf{$w_{k,1}$} & Initial lag-0 weight & Normal(0, 0.5) & Immediate effect of covariate \textit{k} & Starting point for random walk; regularizes initial lag effect. \\
\addlinespace
\textbf{$w_{k,l}$} & Subsequent lag weights & Random Walk: Normal($w_{k,l-1}$, $\sigma_w[k]$) & Effect of covariate \textit{k} at lag \textit{l} & Enforces smoothness across lags via random walk prior; prevents extreme swings in lag structure. \\
\addlinespace
\textbf{$\sigma_w[k]$} & Random walk SD for lag structure & Exponential(2) & Smoothness of lag trajectory & Shrinks toward smooth lag structure; prevents overfitting across lags. \\
\addlinespace
\textbf{$w_{unlagged}$} & Unlagged covariate weights & Normal(0, 0.5) & Effect of static spatial covariates & Regularizes effects of non-temporal covariates (e.g., urban form, land cover). \\
\bottomrule
\end{tabular}
\caption{Prior distributions for model parameters.}
\label{tab:priors}
\end{table}


\subsection{Exploratory Analysis}

Prior to modeling, we compute correlations between each variable and mosquito data, and examine correlations among all variables. When two variables exhibit strong correlation, we may optionally remove one to avoid collinearity issues.

\subsection{Model Fitting}

The Bayesian model is implemented in Stan and fitted via the \texttt{cmdstanr} R interface using MCMC (NUTS sampler). In the current implementation, $\delta_0$ and $\delta_1$ are estimated jointly with all other parameters, so no fixing is required. MCMC settings: 2 chains, 1000 warmup iterations, 1000 sampling iterations, \texttt{adapt\_delta = 0.95}, \texttt{max\_treedepth = 12}. Two model variants are available: with or without the temporal random effect $v_t$ (\texttt{use\_temporal\_re}); the current default omits it.

\subsection{Additional Considerations}

When including lagged coefficients for the same variables, we must correct for temporal association. One option is to constrain the trajectory of coefficients (for a given environmental variable) using a polynomial function; see the ARBOALVO publication for this approach. Additionally, we may consider a hurdle Poisson model, which separately models the probability of zero and the distribution of non-zero values. Such models are well-suited when data contain excess zeros not sufficiently captured by standard models.

% \section{Alternative Approach: Household-Level Model}

% \subsection{Motivation}

% Since data generation starts at the household level, we can model the process mechanistically at that level and aggregate upward.

% \subsection{Latent State Variables}

% We define two latent state variables at the household level: $z_{ib,t} \in \{0,1\}$ indicating whether household $i$ in block $b$ has mosquito larvae at time $t$, and $v_{ib,t} \in \{0,1\}$ indicating whether the household was visited between time $t-1$ and $t$. The application of active control (removal) is defined deterministically as:
% \begin{equation}
% c_{ib,t} = v_{ib,t} \cdot z_{ib,t}.
% \end{equation}

% \subsection{Ecological Transitions}

% The probability of colonization (larvae appearing in a previously uninfested household) and persistence (larvae remaining despite control) are given by:
% \begin{equation}
% P(z_{ib,t} = 1 \mid z_{ib,t-1} = 0) = \text{Logit}^{-1}(\gamma_{bt})
% \end{equation}
% \begin{equation}
% P(z_{ib,t} = 1 \mid z_{ib,t-1} = 1) = \text{Logit}^{-1}(\phi_{bt} - \theta c_{ib,t-1}).
% \end{equation}
% The state $z_{ib,t}$ advances as:
% \begin{equation}
% z_{ib,t} \mid z_{ib,t-1}, v_{ib,t-1} \sim \text{Bernoulli}\left(\text{Logit}^{-1}(\phi_{bt} - \theta c_{ib,t-1})z_{ib,t-1} + \text{Logit}^{-1}(\gamma_{bt})(1 - z_{ib,t-1})\right).
% \end{equation}

% \subsection{Visit Process}

% The probability of a new visit (systematic + reactive) and revisit are given by:
% \begin{equation}
% P(v_{ib,t} = 1 \mid v_{ib,t-1} = 0) = \text{Logit}^{-1}(\alpha_0 + \alpha_1 C_{bt})
% \end{equation}
% \begin{equation}
% P(v_{ib,t} = 1 \mid v_{ib,t-1} = 1) = \text{Logit}^{-1}(\alpha_0 + \alpha_1 C_{bt} + \alpha_2 v_{ib,t-1}),
% \end{equation}
% which yields:
% \begin{equation}
% v_{ib,t} \mid v_{ib,t-1} \sim \text{Bernoulli}\left(\text{Logit}^{-1}(\alpha_0 + \alpha_1 C_{bt} + \alpha_2 v_{ib,t-1})\right).
% \end{equation}

% \subsection{Environmental Drivers}

% Larvae colonization and persistence are modeled as purely driven by ecological covariates with spatio-temporal smoothing:
% \begin{align}
% \gamma_{bt} &= \alpha_0^\gamma + \mathbf{x}_{bt}^T \boldsymbol{\alpha}_{\gamma\phi} + s_b + u_t \\
% \phi_{bt} &= \alpha_0^\phi + \mathbf{x}_{bt}^T \boldsymbol{\alpha}_{\gamma\phi} + s_b + u_t
% \end{align}
% Note that we use different intercepts (persistence easier than colonization?) but shared ecological drivers and shared random effects as a deliberate design choice.

% \subsection{Distributed Lag Covariates (Household Model)}

% As in the original model, we separate covariates that require temporal lagging (e.g., rainfall, temperature) from those that do not (e.g., static land cover, urban form). Let $K_L$ index lagged covariates and $K_U$ index unlagged covariates. Then:
% \begin{align}
% \gamma_{bt} &= \alpha_0^\gamma + \sum_{k\in K_L}\sum_{\ell=0}^{L} w_{k\ell}^{\gamma}\,X_{k,b,t-\ell} + \sum_{k\in K_U} w_{k}^{\gamma}\,X_{k,b} + s_b + u_t, \\
% \phi_{bt} &= \alpha_0^\phi + \sum_{k\in K_L}\sum_{\ell=0}^{L} w_{k\ell}^{\phi}\,X_{k,b,t-\ell} + \sum_{k\in K_U} w_{k}^{\phi}\,X_{k,b} + s_b + u_t.
% \end{align}
% This mirrors the original distributed lag model, but is applied separately to colonization and persistence. Lag weights $w_{k\ell}$ can be regularized (e.g., polynomial or smoothness constraints) to avoid identifiability issues across lags.

% \subsection{Observation and Aggregation}

% We define an observation variable at the household level:
% \begin{equation}
% y_{ib,t} = \mathbb{1}\{z_{ib,t} = 1, v_{ib,t} = 1\}.
% \end{equation}
% We observe the sum over households per block:
% \begin{equation}
% Y_{bt}^* = \sum_{i=1}^{N_b} y_{ib,t}.
% \end{equation}
% The latent variable of interest is the fraction of households with larvae (regardless of observation):
% \begin{equation}
% p_{bt} = \frac{\sum_{i=1}^{N_b} z_{ib,t}}{N_b}.
% \end{equation}

% \subsection{Implementation Challenge}

% Either build a POMP model (but this may be infeasible for hundreds of blocks with hundreds of households) or integrate over households analytically:
% \begin{equation}
% P(z_{ib,t} = 1, v_{ib,t} = 1) = P(z_{ib,t} = 1) \cdot P(v_{ib,t} = 1 \mid z_{ib,t} = 1).
% \end{equation}

% \subsection{Take Expectations per Block (Mean-Field Closure)}

% Define the block-level expectations over households:
% \begin{align}
% p_{bt} &= \mathbb{E}[z_{ib,t}] = \Pr(z_{ib,t}=1) = \frac{1}{N_b}\sum_{i=1}^{N_b} z_{ib,t}, \\
% q_{bt} &= \mathbb{E}[v_{ib,t}] = \Pr(v_{ib,t}=1) = \frac{1}{N_b}\sum_{i=1}^{N_b} v_{ib,t}, \\
% r_{bt} &= \mathbb{E}[c_{ib,t}] = \Pr(c_{ib,t}=1) = \mathbb{E}[z_{ib,t}v_{ib,t}] = \frac{1}{N_b}\sum_{i=1}^{N_b} z_{ib,t}v_{ib,t}.
% \end{align}
% Since $r_{bt}$ is the expected value of a product of two Bernoulli variables, we need a closure:
% \begin{itemize}
% \item \textbf{Independence:} $r_{bt} \approx p_{bt}q_{bt}$.
% \item \textbf{Dependent (approx.):} $r_{bt} \approx p_{bt}q_{bt} + \rho\,p_{bt}(1-p_{bt})q_{bt}(1-q_{bt})$, using $\mathrm{Cov}(X,Y)=\rho\,\mathrm{Var}(X)\,\mathrm{Var}(Y)$ and $\mathrm{Var}(X)=p(1-p)$.
% \end{itemize}

% \subsection{Expectations for Visit Dynamics}

% By the law of total expectation,
% \begin{align}
% q_{bt} &= \mathbb{E}[v_{ib,t}] = \sum_{y\in\{0,1\}} \mathbb{E}[v_{ib,t}\mid v_{ib,t-1}=y]\,\Pr(v_{ib,t-1}=y) \\
% &= \text{Logit}^{-1}(\alpha_0+\alpha_1 C_{bt})\,(1-q_{bt-1}) + \text{Logit}^{-1}(\alpha_0+\alpha_1 C_{bt}+\alpha_2)\,q_{bt-1}.
% \end{align}
% This is a weighted average of transition probabilities for previously unvisited vs. previously visited households.

% \subsection{Expectations for Larvae Dynamics}

% Similarly,
% \begin{align}
% p_{bt} &= \mathbb{E}[z_{ib,t}] = \text{Logit}^{-1}(\gamma_{bt})\,(1-p_{bt-1}) 
%  + \mathbb{E}\big[\text{Logit}^{-1}(\phi_{bt}-\theta v_{ib,t-1}) \mid z_{ib,t-1}=1\big] p_{bt-1}.
% \end{align}
% Since $v_{ib,t-1}\in\{0,1\}$ and $v$ does not explicitly depend on $z$, we can split:
% \begin{equation}
% \mathbb{E}\big[\text{Logit}^{-1}(\phi_{bt}-\theta v_{ib,t-1})\big] = (1-q_{bt-1})\,\text{Logit}^{-1}(\phi_{bt}) + q_{bt-1}\,\text{Logit}^{-1}(\phi_{bt}-\theta).
% \end{equation}
% Therefore,
% \begin{align}
% p_{bt} &= \text{Logit}^{-1}(\gamma_{bt})\,(1-p_{bt-1}) \\
% &\quad + p_{bt-1}\Big[(1-q_{bt-1})\,\text{Logit}^{-1}(\phi_{bt}) + q_{bt-1}\,\text{Logit}^{-1}(\phi_{bt}-\theta)\Big].
% \end{align}

% \subsection{Observation and Simulation Equations}

% Under the independence closure $r_{bt}=p_{bt}q_{bt}$:
% \begin{align}
% q_{bt} &= (1-q_{bt-1})\,\text{Logit}^{-1}(\alpha_0+\alpha_1 C_{bt}) + q_{bt-1}\,\text{Logit}^{-1}(\alpha_0+\alpha_1 C_{bt}+\alpha_2), \\
% p_{bt} &= \text{Logit}^{-1}(\gamma_{bt})\,(1-p_{bt-1}) + p_{bt-1}\Big[(1-q_{bt-1})\,\text{Logit}^{-1}(\phi_{bt}) + q_{bt-1}\,\text{Logit}^{-1}(\phi_{bt}-\theta)\Big], \\
% r_{bt} &= p_{bt}q_{bt}, \\
% \gamma_{bt} &= \alpha_0^\gamma + \mathbf{x}_{bt}^T \boldsymbol{\alpha}_{\gamma\phi} + s_b + u_t, \\
% \phi_{bt} &= \alpha_0^\phi + \mathbf{x}_{bt}^T \boldsymbol{\alpha}_{\gamma\phi} + s_b + u_t.
% \end{align}
% The observation model is then
% \begin{equation}
% Y_{bt} \sim \text{Binomial}(N_b, r_{bt}).
% \end{equation}
% \textit{Note:} We do not observe the number of visits per block, so the Markov chain for $q_{bt}$ may be weakly identifiable.

% \subsection{Special Case: No Revisiting}

% If revisits are negligible ($\alpha_2=0$ and no dependence on $v_{ib,t-1}$), then
% \begin{align}
% q_{bt} &= \text{Logit}^{-1}(\alpha_0+\alpha_1 C_{bt}), \\
% p_{bt} &= \text{Logit}^{-1}(\gamma_{bt})\,(1-p_{bt-1}) + p_{bt-1}\Big[(1-q_{bt-1})\,\text{Logit}^{-1}(\phi_{bt}) + q_{bt-1}\,\text{Logit}^{-1}(\phi_{bt}-\theta)\Big], \\
% r_{bt} &= p_{bt}q_{bt}.
% \end{align}

% \subsection{Model Assumptions}

% \begin{itemize}
% \item Households are interchangeable within a block.
% \item Visits are exogenous (scheduled or externally driven); detection is perfect if visited.
% \item Binomial thinning: each household can be visited at most once per month.
% \item Within a month, $z_{ib,t}$ and $v_{ib,t}$ are independent (otherwise use a bivariate Bernoulli).
% \item Between months, $z_{ib,t}$ depends on $v_{ib,t-1}$ through control, but $v_{ib,t}$ does not depend on $z_{ib,t-1}$. In reality, larvae $\rightarrow$ dengue $\rightarrow$ higher visit probability, so this assumption may be violated.
% \end{itemize}

% \subsection{Model Comparison}

% \begin{table}[h]
% \centering
% \small
% \begin{tabular}{@{}p{2.5cm}p{5.5cm}p{5.5cm}@{}}
% \toprule
% \textbf{Aspect} & \textbf{Original Model} & \textbf{Tijs' New Model} \\ 
% \midrule
% Level of modelling & Block-time, households integrated out implicitly & Household-level, explicit aggregation to blocks \\
% \midrule
% Latent state & $p_{bt}$ = Pr(larvae present at household) & 1. Ecological: $z_{ib,t} \in \{0,1\}$ (larvae present)\newline 2. Visit: $v_{ib,t} \in \{0,1\}$ (household visit) \\
% \midrule
% Dynamics & No explicit colonization/extinction\newline Temporal structure via random effects\newline Ecology drives $p_{bt}$ & Colonization: $P(z_{ib,t}=1|z_{ib,t-1}=0)$\newline Persistence with control: $P(z_{ib,t}=1|z_{ib,t-1}=1)$\newline Visit probability depends on dengue burden and previous visits \\
% \midrule
% Observation & Observation bias via systematic/reactive inspections\newline Dengue cases affect observation, not ecology & Define: $y_{ib,t} = \mathbb{1}\{z_{ib,t}=1, v_{ib,t}=1\}$\newline Observe: $Y_{bt}^* = \sum y_{ib,t}$\newline True hidden Markov model + partial observation \\
% \midrule
% Benefits & 1. Computation manageable\newline 2. Identifiability issues avoided\newline 3. Best for estimating $p_{bt}$ robustly\newline 4. Scaling\newline 5. Inference (not simulation) & 1. Understanding revisits \& removal\newline 2. Understanding control effectiveness in entomology \\
% \midrule
% Risks & 1. No distinction in why $p_{bt}$ changes\newline 2. Intervention effects absorbed into observation bias\newline 3. Need good non-dengue period & 1. Computation: $\sim$2,000,000 latent states\newline 2. Identifiability: never observe $z_{ib,t}=0, v_{ib,t}=1$ or $z_{ib,t}=1, v_{ib,t}=0$ \\
% \bottomrule
% \end{tabular}
% \caption{Comparison of modeling approaches}
% \end{table}

\section{Discussion and Limitations}

Both modeling approaches have important considerations. The block-level model is computationally tractable and avoids certain identifiability issues by anchoring inference during non-outbreak periods and using effort ratios instead of estimating visit probabilities directly. However, it does not distinguish mechanistically why $p_{bt}$ changes over time, and intervention effects are absorbed into the observation bias term. The household-level model offers greater mechanistic detail but faces computational challenges (approximately 2 million latent states for 1,410 blocks with 40 houses each over 36 months) and potential identifiability issues since we never observe households that are visited but have no larvae, or vice versa.

Several limitations apply to the block-level approach. First, we assume that all households are visited monthly through systematic surveillance, but this may not hold in practice. Second, we do not correct for the possibility that multiple dengue cases $C_{bt}$ might originate from the same household, which would violate our assumption that each case triggers inspection of a distinct household. These limitations suggest that the model's correction for reactive surveillance bias may be imperfect, particularly in settings with irregular coverage or spatial clustering of cases.

A key question remains: are we overdoing it by modeling household persistence, visit timing, and explicit colonization/extinction with the mechanistic model? The answer depends on the inferential goals. If the primary objective is robust estimation of $p_{bt}$ for prediction and surveillance evaluation, the simpler block-level model may suffice. If the goal is to understand the mechanistic drivers of mosquito dynamics and evaluate control effectiveness, the added complexity of the household-level model may be justified despite its computational cost.



\part*{Part~II: Frequentist GLMM --- Prior Elicitation}
\addcontentsline{toc}{part}{Part~II: Frequentist GLMM}

\section{Role of the GLMM: Bridging Frequentist and Bayesian Inference}

The Bayesian hierarchical model in Part~I requires prior distributions for covariate effects ($\beta$, $w_{k,\ell}$), surveillance bias parameters ($\delta_0$, $\delta_1$), and variance components ($\sigma_u$, $\sigma_v$). While domain knowledge motivates the general shape of these priors (see Table~\ref{tab:priors}), their location and scale benefit from data-driven calibration.

The frequentist GLMM described in this part serves as a fast, computationally efficient tool to:
\begin{enumerate}
\item \textbf{Screen covariates}: identify which environmental predictors and lags have the strongest associations with mosquito presence, informing variable selection for the Bayesian model.
\item \textbf{Elicit prior locations}: fixed-effect estimates from the GLMM (on the log-odds scale) provide empirically grounded starting points for prior means in the Bayesian model. Because the GLMM uses maximum likelihood without shrinkage, its estimates can be used to centre priors rather than adopted directly as posterior summaries.
\item \textbf{Calibrate variance components}: GLMM estimates of $\sigma_u$ (block heterogeneity) and temporal autocorrelation inform the scale hyperpriors \texttt{Exponential(2)} and \texttt{Exponential(3)} used in Stan.
\end{enumerate}

In short, the GLMM is run first as an exploratory and prior-elicitation step; its outputs then feed into the Bayesian model rather than being the final inferential target.

\section{GLMM for Mosquito Immature Stages Presence}

This section describes the maximum-likelihood GLMM implemented in \texttt{GLMM.r}. The model is estimated at block--month level using \texttt{glmmTMB} in R.

\subsection{Data Structure and Covariates}

For each observation (block $b$, month $t$), the model uses:
\begin{itemize}
\item $y_{bt}$: positive houses (\texttt{Houses\_pos\_IS}),
\item $n_{bt} = N_b^{HH} + \kappa C_{bt}$: total inspection events, where $N_b^{HH}$ is the number of inspected houses (\texttt{Inspected\_houses}), $C_{bt}$ are dengue cases, and $\kappa$ is a scaling factor (see Section~\ref{sec:expansion}),
\item \textbf{lagged environmental covariates} (lags $\ell = 0, 1, 2$, i.e.\ $L = 2$):
  \begin{itemize}
    \item \texttt{rainy\_days\_cat}: number of rainy days per month, entered as an ordered categorical factor (not standardized),
    \item \texttt{avg\_VPD}: average vapour pressure deficit (continuous, z-standardized before lagging),
    \item \texttt{precip\_max\_day}: maximum daily precipitation in the month (continuous, z-standardized before lagging),
  \end{itemize}
\item \textbf{unlagged covariates}:
  \begin{itemize}
    \item \texttt{is\_urban}: binary indicator of urban block (not standardized),
    \item \texttt{water\_containers}: count of water containers (continuous, z-standardized),
  \end{itemize}
\item \texttt{reactive\_shift}: outbreak-intensity covariate added to reactive rows (see Section~\ref{sec:expansion}).
\end{itemize}

Only continuous numeric variables (\texttt{avg\_VPD}, \texttt{precip\_max\_day}, \texttt{water\_containers}) are z-score standardized prior to lag creation.
Categorical variables (\texttt{rainy\_days\_cat}) and binary indicators (\texttt{is\_urban}) are not standardized.
Fixed-effect coefficients for z-standardized predictors are interpreted on a one-standard-deviation scale (change in log-odds for a $+1$ SD increase).

\subsection{Data Expansion: Baseline and Reactive Splits}
\label{sec:expansion}

Rather than fitting the surveillance mixture directly, the code expands each original (block, month) observation into separate baseline and reactive rows based on the reactive fraction $\omega_{bt}$:
\begin{equation}
\omega_{bt} = \min\!\left(\frac{\kappa\, C_{bt}}{N_b^{HH} + \kappa\, C_{bt}},\; 1\right),
\end{equation}
which is the proportion of total inspection events attributable to reactive (case-triggered) surveillance.

\noindent\textbf{Baseline row (systematic surveillance):}
\begin{equation}
n_{bt}^{(base)} = \lfloor (1 - \omega_{bt})\, n_{bt} \rfloor, \quad y_{bt}^{(base)} = \lfloor y_{bt}(1-\omega_{bt}) \rfloor, \quad \texttt{reactive\_shift} = 0.
\end{equation}

\noindent\textbf{Reactive row (case-triggered inspections, created only when $C_{bt} > 0$):}
\begin{equation}
n_{bt}^{(react)} = \lfloor \omega_{bt}\, n_{bt} \rfloor, \quad y_{bt}^{(react)} = \lfloor y_{bt}\, \omega_{bt} \rfloor, \quad \texttt{reactive\_shift} = \log(1 + C_{bt}).
\end{equation}

The \texttt{reactive\_shift} covariate enters the linear predictor only for reactive rows, allowing the model to estimate the systematic log-odds difference between surveillance regimes.
Rows with $n_{bt}^{(base)} = 0$ or $n_{bt}^{(react)} = 0$ (after flooring) are dropped, given that it's implausible to find mosquitoes when there are no houses to survey. 

\subsubsection{Example Configuration: $\kappa = 2$}

With $\kappa = 2$, each dengue case is assumed to trigger on average 2 additional household inspections.
The total trial count becomes $n_{bt} = N_b^{HH} + 2C_{bt}$, and the reactive fraction is:
\begin{equation}
\omega_{bt} = \frac{2\,C_{bt}}{N_b^{HH} + 2\,C_{bt}}.
\end{equation}

For months with $C_{bt} > 0$, each (block, month) observation is split into two rows.
For months with $C_{bt} = 0$, $\omega_{bt} = 0$ and only a baseline row is produced (the expansion is inactive).

\noindent\textbf{Concrete example:} suppose $N_b^{HH} = 100$ inspected houses and $C_{bt} = 5$ cases. Then:
\begin{align*}
n_{bt} &= 100 + 2 \times 5 = 110, \\
\omega_{bt} &= 10/110 \approx 0.091, \\
n_{bt}^{(base)} &= \lfloor 0.909 \times 110 \rfloor = 100, \quad y_{bt}^{(base)} = \lfloor y_{bt} \times 0.909 \rfloor, \\
n_{bt}^{(react)} &= \lfloor 0.091 \times 110 \rfloor = 10, \quad y_{bt}^{(react)} = \lfloor y_{bt} \times 0.091 \rfloor, \quad \texttt{reactive\_shift} = \log(6) \approx 1.79.
\end{align*}
The baseline row retains the original household universe ($N_b^{HH}$) as its trial count, and the 10 reactive inspections form a separate row with a non-zero \texttt{reactive\_shift}.
The coefficient $\delta$ on \texttt{reactive\_shift} then estimates the log-odds difference between reactive and systematic surveillance, conditional on all other covariates.

\subsubsection{Surveillance-Mixture Representation (when $\kappa > 0$)}

When $\kappa > 0$, the data expansion encodes a \textbf{latent two-process data-generating mechanism}: each observation with $C_{bt} > 0$ is decomposed into a baseline component (routine inspections) and a reactive component (case-triggered inspections), both modeled as separate binomial contributions within the same GLMM.
After prediction, the model yields two fitted probabilities per original (block, month) unit:
\begin{itemize}
\item $\hat{p}_{bt}$: baseline (ecological) fitted probability,
\item $\hat{p}_{bt}^R$: reactive fitted probability (inflated by targeting bias).
\end{itemize}

\subsection{Critical Assumption: Proportional Allocation of Positives}

The data expansion assumes that observed positives are split proportionately to the reactive fraction:
\begin{equation}
\frac{y_{bt}^{(base)}}{n_{bt}^{(base)}} = \frac{y_{bt}^{(react)}}{n_{bt}^{(react)}} = \frac{y_{bt}}{n_{bt}}.
\end{equation}

This is a \textbf{strong and problematic} assumption because:

\begin{enumerate}
\item \textbf{Reactive inspections may be targeted.} The model assumes each case originates from a distinct household, so reactive inspections are spread across $\kappa C_{bt}$ different locations within the block. If cases instead cluster in the same household or a small set of high-risk households --- which is plausible given the spatial aggregation of dengue transmission --- reactive inspections would be concentrated there and their positivity rate could be \textit{higher} than the block-wide average assumed by proportional allocation. The extent of this bias depends on how spatially concentrated cases are within a block.

\item \textbf{Non-independence of effort and detection.} The number of reactive inspections is driven by observed cases $C_{bt}$, which are themselves outputs of the underlying positivity $p_{bt}$. This confounds effort allocation with the ecological signal.

\item \textbf{Distortion of the reactive shift.} Proportional allocation artificially constrains $p_{bt}^R = p_{bt}$, so the \texttt{reactive\_shift} coefficient $\delta$ captures only the residual difference after this artificial equalization. The magnitude and sign of $\delta$ may therefore be unreliable.
\end{enumerate}

For now, this limitation should be acknowledged when interpreting the reactive shift. Future work should prioritize obtaining data that labels inspections by type (systematic vs.\ reactive) to enable more rigorous estimation.

\subsection{Linear Predictor and GLMM Model}

For each row $i$ (baseline or reactive), the linear predictor is:
\begin{align}
\eta_i &= \alpha
  + \sum_{k \in \mathcal{L}} \sum_{\ell=0}^{L} \beta_{k\ell}\, x_{k,b(i),t(i)-\ell}
  + \sum_{k \in \mathcal{U}} \beta_k\, x_{k,b(i)}
  + \delta \cdot \text{reactive\_shift}_i \notag \\
&\quad + v_{t(i)} + \phi_{b(i),t(i)},
\end{align}
where:
\begin{itemize}
\item $\mathcal{L} = \{\texttt{rainy\_days\_cat},\; \texttt{avg\_VPD},\; \texttt{precip\_max\_day}\}$ are the lagged covariates (lags $\ell = 0, 1, 2$),
\item $\mathcal{U} = \{\texttt{is\_urban},\; \texttt{water\_containers}\}$ are the unlagged covariates,
\item $\delta$ = coefficient for \texttt{reactive\_shift} $= \log(1 + C_{bt})$,
\item $v_t \sim \text{Normal}(0, \sigma_v)$ = i.i.d.\ random intercept for month $t$ (\texttt{year\_month}), capturing shared temporal shocks not explained by covariates,
\item $\phi_{b,t}$ = AR(1) temporal random effect \emph{within} block $b$, modeled as:
  \begin{equation}
    \phi_{b,t} = \rho\, \phi_{b,t-1} + \varepsilon_{b,t}, \quad \varepsilon_{b,t} \sim \text{Normal}(0, \sigma_\phi),
  \end{equation}
  where $\rho \in (-1, 1)$ is the autoregressive parameter estimated jointly for all blocks. This is specified in \texttt{glmmTMB} as \texttt{ar1(year\_month\_ar1 + 0 | block)}: the \texttt{+ 0} suppresses a redundant random intercept within the AR(1) term (the global temporal mean is already captured by $v_t$).
\end{itemize}

\subsubsection{Random Effects: Structure and Options}

The model supports four optional random effects components, each targeting a different source of dependency in the data.
Table~\ref{tab:re_options} summarizes them; the current defaults are marked.

\begin{table}[h]
\centering
\small
\begin{tabular}{p{3.5cm} p{2cm} p{3cm} p{5.5cm}}
\toprule
\textbf{Component} & \textbf{Default} & \textbf{glmmTMB term} & \textbf{Correlation structure} \\
\midrule
Block random intercept & \textbf{OFF} & \texttt{(1 | block)} & i.i.d.\ Normal; captures stable between-block differences in baseline positivity \\
\addlinespace
Month random intercept & \textbf{ON} & \texttt{(1 | year\_month)} & i.i.d.\ Normal; captures shared temporal shocks (e.g.\ city-wide weather anomalies) not explained by fixed covariates \\
\addlinespace
Temporal AR(1) & \textbf{ON} & \texttt{ar1(year\_month\_ar1 + 0 | group)} & First-order autoregression in time within each group; $\text{Corr}(\phi_{b,t},\phi_{b,t+k}) = \rho^k$. The \texttt{+ 0} suppresses a redundant intercept (already captured by the month RE above). Group is either \texttt{block} (separate AR(1) per block, same $\rho$) or \texttt{"all"} (single global AR(1)) \\
\addlinespace
Spatial AR & \textbf{OFF} & \texttt{exp(xy + 0 | spatial)} & Exponential decay with distance: $\text{Corr}(s_b, s_{b'}) = \exp(-d_{bb'}/r)$, where $d_{bb'}$ is Euclidean distance (m) and $r$ is the estimated range. Requires projected coordinates from shapefile. Alternative: Mat\'{e}rn covariance (\texttt{mat(\ldots)}) \\
\bottomrule
\end{tabular}
\caption{Available random effects components and their correlation structures. Current defaults shown in the Default column.}
\label{tab:re_options}
\end{table}

The general linear predictor including all four components is:
\begin{equation}
\eta_i = \alpha + \sum_{k \in \mathcal{L}} \sum_{\ell=0}^{L} \beta_{k\ell}\, x_{k,b,t-\ell} + \sum_{k \in \mathcal{U}} \beta_k x_{k,b} + \delta \cdot \text{reactive\_shift}_i + u_b + v_t + \phi_{b,t} + \psi_b,
\end{equation}
where $u_b \sim \mathcal{N}(0,\sigma_u)$ (block intercept), $v_t \sim \mathcal{N}(0,\sigma_v)$ (month intercept), $\phi_{b,t}$ (within-block AR(1)), and $\psi_b$ (spatial effect) are each optional.
In the current configuration only $v_t$ and $\phi_{b,t}$ are active.

In R/\texttt{glmmTMB} notation, the current formula is:
\begin{verbatim}
cbind(y_bt, n_trials - y_bt) ~
  rainy_days_cat_lag0 + rainy_days_cat_lag1 + rainy_days_cat_lag2 +
  avg_VPD_lag0 + avg_VPD_lag1 + avg_VPD_lag2 +
  precip_max_day_lag0 + precip_max_day_lag1 + precip_max_day_lag2 +
  is_urban + water_containers + reactive_shift +
  (1 | year_month) +
  ar1(year_month_ar1 + 0 | block)
\end{verbatim}

\subsection{Observation Model}

The likelihood for each row $i$ is:
\begin{equation}
y_i \sim \text{Binomial}(n_i,\; \pi_i), \quad \text{where} \quad \pi_i = \text{logit}^{-1}(\eta_i).
\end{equation}
In \texttt{glmmTMB}, this is specified as \texttt{cbind(y\_i, n\_i - y\_i) \textasciitilde{} \ldots}, passing successes (positive houses) and failures (negative houses) separately.
The response therefore weights observations by trial count: a block with 5/50 positive houses contributes more information than one with 5/10, even though the observed proportion is the same.

\subsection{Estimation via glmmTMB}

The model is fit using:
\begin{itemize}
\item \textbf{Optimizer:} \texttt{nlminb} (via TMB backend), with \texttt{iter.max = eval.max = 10\,000},
\item \textbf{Likelihood:} Binomial with configurable link function (default: \texttt{logit}),
\item \textbf{Random effects:} Month i.i.d.\ intercept $v_t$ and within-block AR(1) $\phi_{b,t}$ (both currently enabled); block random intercept currently disabled,
\item \textbf{Regularization:} None explicit; standard frequentist inference with MLE via the Laplace approximation in TMB.
\end{itemize}

\subsection{Model Outputs}

The script produces:
\begin{itemize}
\item Fitted ecological probabilities $\hat{p}_{bt}$ (from baseline rows),
\item Fitted reactive probabilities $\hat{p}_{bt}^R$ (from reactive rows, when $\kappa > 0$),
\item A weighted-average fitted probability $\hat{\pi}_{bt} = (1 - \omega_{bt})\hat{p}_{bt} + \omega_{bt}\hat{p}_{bt}^R$ at the original (block, month) level,
\item Fixed-effect coefficient table with odds ratios,
\item Diagnostic plots: time-series of fitted vs.\ observed probabilities, random effect estimates, QQ plots.
\end{itemize}

\subsection{Model Flexibility}

All key choices are controlled via the \texttt{cfg} list in \texttt{GLMM.r}.
Current defaults and available options:
\begin{itemize}
\item \texttt{include\_block\_re = FALSE}: toggle block-level i.i.d.\ random intercept $(1 \mid \text{block})$,
\item \texttt{include\_time\_re = TRUE}: toggle month-level i.i.d.\ random intercept $(1 \mid \text{year\_month})$,
\item \texttt{include\_ar1\_temporal = TRUE}: toggle temporal AR(1) within groups; correlation decays as $\rho^k$ over $k$ months,
\item \texttt{ar1\_group = "block"}: grouping for AR(1); \texttt{"block"} = separate AR(1) per block (same $\rho$), \texttt{"global"} = single AR(1) across all observations,
\item \texttt{include\_spatial\_ar = FALSE}: toggle exponential spatial covariance $\exp(-d/r)$ across blocks (requires shapefile with block coordinates); alternative Mat\'{e}rn covariance is available in the code but commented out,
\item \texttt{lag\_vars}: variables for which distributed lags are created (default: \texttt{rainy\_days\_cat}, \texttt{avg\_VPD}, \texttt{precip\_max\_day}),
\item \texttt{unlagged\_vars}: variables entered without lagging (default: \texttt{is\_urban}, \texttt{water\_containers}),
\item \texttt{numeric\_vars}: variables to z-standardize (default: \texttt{precip\_max\_day}, \texttt{avg\_VPD}, \texttt{water\_containers}),
\item \texttt{max\_lag = 2}: maximum lag $L$ for all lagged covariates,
\item \texttt{kappa = 0}: reactive inspection scaling; $\kappa = 0$ disables the baseline/reactive split,
\item \texttt{link\_function = "logit"}: link function for binomial family (\texttt{"logit"}, \texttt{"probit"}, \texttt{"cloglog"}, \texttt{"cauchit"}),
\item \texttt{interactions = NULL}: optional list of two-variable interaction pairs,
\item \texttt{iter\_max, eval\_max}: optimizer convergence criteria (default: $10^4$ each).
\end{itemize}

This approach offers computational efficiency and straightforward interpretation of fixed and random effects under a frequentist MLE framework.

\appendix
\section{Appendix: Example GLMM Results and Interpretation}

\textbf{Note:} The results below are from an \emph{older} model configuration (block RE enabled, time RE enabled, no AR(1), lags 0--1, covariates: temperature, humidity, precipitation, NDVI, relative humidity, windspeed).
They do not reflect the current default configuration described in the main text ($\kappa = 0$, max lag $= 2$, covariates: \texttt{rainy\_days\_cat}, \texttt{avg\_VPD}, \texttt{precip\_max\_day}).
This appendix is retained for reference only and will be updated when new results are available.

\subsection{Fixed Effects (Odds-Ratio Interpretation)}

Table~\ref{tab:glmm_example_fixed} reports all fixed effects from \texttt{glmm\_fixed\_effects\_OR\_20260309\_blockRE\_timeRE.csv}, including estimate, standard error, p-value, and odds ratio.

\begin{table}[h]
\centering
\small
\resizebox{\textwidth}{!}{%
\begin{tabular}{lcccc}
\toprule
Term & Estimate ($\beta$) & Std. Error & p-value & OR ($e^{\beta}$) \\
\midrule
(Intercept) & -6.131 & 0.082 & $< 10^{-15}$ & 0.002 \\
\texttt{reactive\_shift} & -0.813 & 0.216 & $1.62 \times 10^{-4}$ & 0.443 \\
\texttt{is\_WUITRUE} & -0.496 & 0.102 & $1.25 \times 10^{-6}$ & 0.609 \\
\texttt{is\_urbanTRUE} & 0.484 & 0.068 & $1.36 \times 10^{-12}$ & 1.623 \\
\texttt{mean\_ndvi\_lag1} & -0.201 & 0.042 & $1.46 \times 10^{-6}$ & 0.818 \\
\texttt{total\_precip\_lag0} & -0.185 & 0.104 & 0.0744 & 0.831 \\
\texttt{avg\_temp\_lag1} & 0.180 & 0.121 & 0.1347 & 1.198 \\
\texttt{has\_aljibesTRUE} & -0.146 & 0.413 & 0.7243 & 0.864 \\
\texttt{rel\_hum\_lag1} & 0.145 & 0.123 & 0.2356 & 1.157 \\
\texttt{rel\_hum\_lag0} & 0.137 & 0.123 & 0.2652 & 1.146 \\
\texttt{water\_shortageTRUE} & 0.126 & 0.099 & 0.2033 & 1.134 \\
\texttt{is\_WITRUE} & -0.111 & 0.064 & 0.0824 & 0.895 \\
\texttt{avg\_temp\_lag0} & 0.107 & 0.137 & 0.4346 & 1.113 \\
\texttt{total\_precip\_lag1} & -0.101 & 0.100 & 0.3135 & 0.904 \\
\texttt{WS2M} & 0.095 & 0.095 & 0.3179 & 1.099 \\
\texttt{mean\_ndvi\_lag0} & -0.070 & 0.044 & 0.1167 & 0.933 \\
\texttt{nr\_aljibes} & 0.024 & 0.049 & 0.6238 & 1.025 \\
\texttt{water\_containers} & 0.006 & 0.027 & 0.8162 & 1.006 \\
\bottomrule
\end{tabular}
}
\caption{All fixed effects from the fitted GLMM (block and time random intercepts included).}
\label{tab:glmm_example_fixed}
\end{table}

\textbf{Odds-ratio level interpretation:} For each term, $\exp(\beta)$ is the multiplicative change in odds for a one-unit increase in the predictor.

\textbf{Standardized-estimate interpretation:} For continuous predictors (NDVI, temperature, precipitation, humidity, windspeed, number of aljibes, number of depositos), one unit corresponds to \textbf{+1 SD} because variables were z-standardized before fitting. Therefore, $\beta=-0.201$ for \texttt{mean\_ndvi\_lag1} means that a +1 SD NDVI increase changes log-odds by -0.201, equivalent to OR $=0.818$ (18.2\% lower odds). Likewise, $\beta=0.180$ for \texttt{avg\_temp\_lag1} implies OR $=1.198$ (19.8\% higher odds) per +1 SD.

\subsection{Baseline vs Reactive Fitted Probabilities}

For each original block--month unit with $C_{bt}>0$, the expanded model yields two fitted probabilities:
\begin{itemize}
\item baseline fitted probability $\hat{p}_{bt}$,
\item reactive fitted probability $\hat{p}_{bt}^{R}$.
\end{itemize}


\end{document}
